% ========================================================================
% ELEMENTOS PRÉ-TEXTUAIS
% Capa e folha de rosto conforme artigo TCC IFPI
% ========================================================================

% ========================================================================
% CAPA (Obrigatória em artigos do IFPI segundo o manual TCC)
% ========================================================================
\imprimircapa

% ========================================================================
% FOLHA DE ROSTO (Obrigatória)
% ========================================================================
\imprimirfolhaderosto*

% ========================================================================
% AGRADECIMENTOS (OPCIONAL - EM ARTIGOS DEVE IR NO FINAL)
% ========================================================================
% O manual do IFPI diz que agradecimentos em artigos vão como elemento pós-textual,
% portanto mude-os para o final do `main.tex` se for usar.

% ========================================================================
% CABEÇALHO DO ARTIGO (Título, autores, e-mails)
% ========================================================================
\imprimircabecalhoartigo

% ========================================================================
% RESUMO EM PORTUGUÊS E PALAVRAS-CHAVE
% ========================================================================
\setlength{\absparsep}{18pt} % Ajusta espaçamento entre parágrafos do resumo

\begin{resumo}
    Este artigo científico tem como objetivo [DESCREVER O OBJETIVO PRINCIPAL DA PESQUISA]. A metodologia utilizada baseia-se em [DESCREVER A METODOLOGIA: pesquisa bibliográfica, estudo de caso, pesquisa experimental, etc.]. Como referencial teórico, serão utilizados os trabalhos de [CITAR PRINCIPAIS AUTORES DA ÁREA]. A pesquisa justifica-se pela [APRESENTAR A JUSTIFICATIVA E RELEVÂNCIA DO TEMA]. Os resultados incluem [DESCREVER OS RESULTADOS ALCANÇADOS]. Este estudo contribui para [APRESENTAR A CONTRIBUIÇÃO PARA A ÁREA DE CONHECIMENTO].
    
    \vspace{\onelineskip}
    \noindent
    \textbf{Palavras-chave}: palavra-chave 1; palavra-chave 2; palavra-chave 3; palavra-chave 4; palavra-chave 5.
\end{resumo}

% ========================================================================
% RESUMO EM INGLÊS (ABSTRACT) E KEYWORDS
% ========================================================================
\begin{resumo}[ABSTRACT]
    \begin{otherlanguage*}{english}
        \setlength{\parindent}{0pt}% Garante que não haverá recuo na primeira linha do abstract
        This scientific article aims to [DESCRIBE THE MAIN OBJECTIVE OF THE RESEARCH]. The methodology used is based on [DESCRIBE THE METHODOLOGY: bibliographic research, case study, experimental research, etc.]. As theoretical framework, the works of [CITE MAIN AUTHORS IN THE FIELD] will be used. The research is justified by [PRESENT THE JUSTIFICATION AND RELEVANCE OF THE THEME]. Results include [DESCRIBE REACHED RESULTS]. This study contributes to [PRESENT THE CONTRIBUTION TO THE KNOWLEDGE AREA].
        
        \vspace{\onelineskip}
        \noindent
        \textbf{Keywords}: keyword 1; keyword 2; keyword 3; keyword 4; keyword 5.
    \end{otherlanguage*}
\end{resumo}

% ========================================================================
% DATA DE APROVAÇÃO
% ========================================================================
\vspace{\onelineskip}
\noindent
Data de aprovação: \@dataaprovacao

\vspace{1cm}

% ========================================================================
% CONFIGURAÇÕES PARA INÍCIO DO TEXTO
% ========================================================================
% Impede quebra de página padrão do abntex2 após pré-textuais
\makeatletter
\renewcommand{\textual}{%
  \pagestyle{ifpihead}
  \aliaspagestyle{chapter}{ifpihead}
  \nouppercaseheads
  \bookmarksetup{startatroot}
}
\makeatother

\textual
